% Irodov user PDF pp.63-64, original printed pp.54-55.
% Final original hydrodynamics section 1.7 problems 335-339.

\exproblem{335}{54}
原书图1.90中，水从宽大水箱 $A$
通过一根截面均匀但有黏性阻力的水平管流出。
沿管道相隔同样距离 $\ell$ 装两根竖直测压管，
其水柱高分别为 $h_2=20\,\mathrm{cm}$、
$h_1=10\,\mathrm{cm}$；
从水箱接口到第一根测压管、从第二根测压管到出口
也各为 $\ell$。水箱自由水面高出管轴
$h_3=35\,\mathrm{cm}$。
求管内水流速率 $v$。
\exsolution
在截面不变的细管内，稳定流动速度 $v$ 不随
轴向位置变化；黏性压头损失与等长管段长度成正比。
两测压管相距 $\ell$，压头损失为
$h_2-h_1=10\,\mathrm{cm}$。
因此沿上游再外推长度 $\ell$，
管入口处的压力水头为
\[
h_{\rm inlet}=h_2+(h_2-h_1)
=2h_2-h_1=30\,\mathrm{cm}.
\]
水箱自由面到管入口还可用于把水从近静止
加速到管速 $v$ 的水头为
\[
\Delta h=h_3-h_{\rm inlet}
=h_3-2h_2+h_1=5\,\mathrm{cm}.
\]
对水箱自由面和管入口应用含上述分配的能量式：
$v^2/(2g)=\Delta h$。
于是
\[
\boxed{v=\sqrt{2g(h_3-2h_2+h_1)}
=\sqrt{2(9.81)(0.05)}
\simeq1.0\,\mathrm{m/s}}.
\]
\excheck 不能把全部 $h_3$ 当作动能水头；
图示测压管读数已显示管内存在黏性沿程损失。

\exproblem{336}{54}
圆形管道的内半径沿轴向按
$r(x)=r_0e^{-\alpha x}$ 缓慢减小，
其中 $\alpha=0.50\,\mathrm{m^{-1}}$。
流体密度、黏度不变，体积流量恒定。
求沿程相隔 $\Delta x=3.2\,\mathrm m$ 的
两个截面的雷诺数比值。
\exsolution
取管内流动的特征长度为管径
$D(x)=2r(x)$，平均速度
$\bar v(x)=Q/[\pi r(x)^2]$，
雷诺数
\[
\mathrm{Re}(x)=\frac{\rho\bar v(x)D(x)}{\eta}
=\frac{2\rho Q}{\pi\eta r(x)}.
\]
因此在均匀流量下，
$\mathrm{Re}(x)\propto1/r(x)\propto e^{\alpha x}$。
两个位置 $x$ 和 $x+\Delta x$ 的比值
\[
\boxed{\frac{\mathrm{Re}(x+\Delta x)}{\mathrm{Re}(x)}
=e^{\alpha\Delta x}=e^{0.5\times3.2}
=e^{1.6}\simeq4.95\approx5}.
\]
\excheck 管径越小，平均速度增长更快于特征长度减小，
故雷诺数增大而非减小；本式以密度、黏度和流量恒定为前提。

\exproblem{337}{54}
半径为 $r_1=1.2\,\mathrm{mm}$ 的小球
在甘油中以速度不超过
$v_1=23\,\mathrm{cm/s}$ 时处于给定的层流条件。
在水中换成半径
$r_2=5.5\,\mathrm{cm}$ 的大球，
问达到相同临界雷诺数时至少需多大速度 $v_2$？
甘油与水黏度分别为
$\eta_1=13.9\,\mathrm P$、
$\eta_2=0.011\,\mathrm P$。
\exsolution
球体绕流的特征长度为直径 $2r$，
雷诺数
$\mathrm{Re}=2\rho vr/\eta$。
按题意取两个液体在\emph{同一临界雷诺数}时比较，
\[
\frac{2\rho_1v_1r_1}{\eta_1}
=\frac{2\rho_2v_2r_2}{\eta_2}.
\]
所以
\[
\boxed{v_2=v_1
\frac{r_1\rho_1\eta_2}{r_2\rho_2\eta_1}}.
\]
代入甘油密度约
$\rho_1=1.26\times10^3\,\mathrm{kg/m^3}$、
水密度 $\rho_2=1.00\times10^3\,\mathrm{kg/m^3}$
及所给黏度，得到
\[
\boxed{v_2\simeq5.0\times10^{-6}\,\mathrm{m/s}
=5\,\mathrm{\mu m/s}}.
\]
\excheck 本题按原书把雷诺数阈值称为
“层流变湍流”，但围绕球的 Stokes 流失效
与充分发展湍流并不必然是同一个转捩点；
计算结果严格属于\emph{雷诺数相似}下的阈值换算。

\exproblem{338}{54}
一颗铅球在动力黏度为
$\eta=13.9\,\mathrm P$ 的甘油中
以终端速率稳定下沉。
定义球直径为雷诺数特征长度，
当临界值 $\mathrm{Re}_{\rm cr}=0.5$ 时
视为超出题定层流范围。
求球仍符合这一范围的最大直径 $d_{\max}$。
\exsolution
记铅密度 $\rho_s$、甘油密度 $\rho_f$。
在 Stokes 阻力近似中，
球重力减浮力与黏滞阻力平衡：
\[
(\rho_s-\rho_f)\frac{\pi d^3}{6}g
=3\pi\eta d\,v_t.
\]
于是
\[
v_t=\frac{(\rho_s-\rho_f)gd^2}{18\eta}.
\]
雷诺数定义为
$\mathrm{Re}=\rho_fv_td/\eta$，
代入后
\[
\mathrm{Re}=
\frac{\rho_f(\rho_s-\rho_f)gd^3}{18\eta^2}.
\]
在 $\mathrm{Re}=\mathrm{Re}_{\rm cr}$ 处，
\[
\boxed{d_{\max}=
\left[\frac{18\mathrm{Re}_{\rm cr}\eta^2}
{\rho_f(\rho_s-\rho_f)g}\right]^{1/3}}.
\]
以 $\rho_s\simeq11.3\times10^3\,\mathrm{kg/m^3}$、
$\rho_f\simeq1.26\times10^3\,\mathrm{kg/m^3}$、
$\eta=1.39\,\mathrm{Pa\,s}$，
\[
\boxed{d_{\max}\simeq5.2\,\mathrm{mm}\approx5\,\mathrm{mm}}.
\]
\excheck 直径因公式中 $d^3$ 出现而为\emph{立方根}；
若误取平方根会破坏量纲。
雷诺数恰为0.5时 Stokes 近似属于临界估算。

\exproblem{339}{54--55}
直径 $d=3.0\,\mathrm{mm}$ 的钢球
由静止开始在橄榄油中下沉，
油的动力黏度为 $\eta=0.90\,\mathrm P$。
按 Stokes 线性阻力模型，
求钢球速率距终端速率相差不到
$n=1.0\%$ 所需时间。
\exsolution
设球体积 $V=\pi d^3/6$、钢密度 $\rho_s$、
油密度 $\rho_f$。
用向下为正写牛顿运动方程：
\[
m\dot v=(\rho_s-\rho_f)Vg-3\pi\eta d\,v,
\qquad m=\rho_sV.
\]
令
$k=3\pi\eta d$、
$\tau=m/k=\rho_s d^2/(18\eta)$。
方程的终端速率为
$v_t=(\rho_s-\rho_f)Vg/k$，
初始 $v(0)=0$ 的解是
\[
v(t)=v_t(1-e^{-t/\tau}).
\]
故相对终端速率差
$(v_t-v)/v_t=e^{-t/\tau}$。
要求它等于 $\varepsilon=0.01$ 时，
\[
\boxed{t=\tau\ln\frac1\varepsilon
=\frac{\rho_s d^2}{18\eta}\ln100}.
\]
取 $\rho_s\simeq7.8\times10^3\,\mathrm{kg/m^3}$、
$\eta=0.090\,\mathrm{Pa\,s}$，
\[
\boxed{t\simeq0.20\,\mathrm s}.
\]
\excheck 模型未计入附加质量和历史阻力。
对题给球径及橄榄油黏度，实际终端雷诺数可能不够低，
因此 Stokes 线性阻力给出的是\emph{题定模型的结果}，
并非无需校准的高精度真实沉降时间。
